\documentclass[11pt]{article}


% Useful packages
\usepackage[utf8]{inputenc}  % para acentos y ñ
\usepackage[english,spanish]{babel}

\usepackage[letterpaper,top=2cm,bottom=2cm,left=2cm,right=2cm,marginparwidth=1.75cm]{geometry}

\usepackage{amsmath}
\usepackage{amsfonts}
\usepackage{amssymb}

\usepackage{listings} % poner codigo
\lstset{
	inputencoding=utf8,
	extendedchars=true,
	literate={á}{{\'a}}1 {é}{{\'e}}1 {í}{{\'i}}1
	{ó}{{\'o}}1 {ú}{{\'u}}1 {ñ}{{\~n}}1
	{Á}{{\'A}}1 {É}{{\'E}}1 {Í}{{\'I}}1
	{Ó}{{\'O}}1 {Ú}{{\'U}}1 {Ñ}{{\~N}}1,
}


\usepackage{float}

\usepackage{graphicx}

%%%%%%%%%%%%% Referencias
\usepackage[colorlinks=true, allcolors=black]{hyperref}
\usepackage{cleveref}

\graphicspath{ {images/} }

\usepackage{geometry}

\usepackage{fancyhdr}
%\setlength{\headheight}{14pt}
\pagestyle{fancy} 

\usepackage{multicol}
\setlength{\columnsep}{0,5cm}

\usepackage{lipsum}

\usepackage[usenames]{color}

%tablas
\usepackage{tabularx} %ancho de tablas
\usepackage{longtable}
\usepackage{ltablex}
\keepXColumns

\usepackage{caption} % títulos de las tablas/figuras
\captionsetup{font=small} % letra de todos los caption en pequeña
\captionsetup[table]{name=Tabla}
\usepackage{subcaption} % este y el anteriori para poner imagenes juntas

%Bibliografia
\usepackage{hyperref}
\usepackage{cite}
%\usepackage{balance}%evita que se partan cosas

%%%%%%%%%%%%%%% NUEVOS COMANDOS
\newcommand{\comillas}[1]{"#1"} %comillas dobles
\newcommand{\mult}[0]{\!\cdot\!} %comillas dobles


%%%%%%%%%%%%%%%%%%%%%%%%%%%%%%%%%%%%%%%%%%%%%%%%%%%%%%%%%%%
%%%%%%%%%%%%%%%%
%%%%%%%%%%%%%%%%%%%%%%%%%%%%%%%%%%%%%%%%%%%%%%%%%%%%%%%%%%%%%%%%%%%%%%%%%%%
\begin{document}
	\renewcommand\contentsname{\huge{Índice General}}
	
	
	
	\begin{center}
		
		
		
		\vspace*{3,5cm} % Espacio entre tope superior y titulo   
		\textbf{\Huge Ecuaciones del movimiento del péndulo simple mediante conservación de la energía y el formalismo lagrangiano}
		
		
		%\bigskip
		%\bigskip
		
		%\textbf{\Huge dominio de la frecuencia}
		
		\bigskip
		\bigskip
		\bigskip
		\bigskip
		\bigskip
		\bigskip
		
		\LARGE{Mario Carracedo Andrés}\\
		\LARGE{}
		
		\bigskip
		\bigskip
		\bigskip
		\bigskip
		
		\includegraphics[height=10cm, scale=2]{imagenes/Pendulo_simple}
		
	\end{center}
	
	
	%%%%%%%%%%%%%%%%%%%%%%%%%%%%%%%%%%%%%%%%%%%%%%%%%%%%%%%%%%%%%%%%%%%%%%%%%%%%
	%%%%%%%%%%%%%%%%%%%%%%%%%%%%%%%%%%%%%%%%%%%%%%%%%%%%%%%%%%%%%%%%%%%%%%%%%%%%
	
	\newpage
	
	\fancyhead{} 
	\fancyhead[l]{Mario Carracedo Andrés}
	\fancyhead[r]{}
	\fancyfoot[c]{\thepage}
	
	% ÍNDICE
	\tableofcontents 
	\newpage
	\section{Consideraciones}
		Para resolver el sistema vamos a tener en cuenta que la cuerda de longitud l, no varía su longitud y que la masa no es un sólido rígido, sino puntual. \\
		\\
		Además tomaremos un sistema de referencia, S(x,y) centrado en el punto de unión de la cuerda con el techo, de la siguiente forma:
		
		\begin{figure}[H]
			\centering
			\includegraphics[scale=0.5]{imagenes/Pendulo_simple_SR}
			\caption{}		
			\label{Fig: Pendulo SR}
		\end{figure}
	\section{Conservación de la energía}
		Lo primero de todo será construir nuestro vector de posición $\vec{r}$, en donde aunque parezca que nuestra masa dependa de $\hat{x}$ y $\hat{y}$, realmente la masa va a describir una cuerva en un plano 2D, por ende quiere decir que vamos a poder tener una coordenada generalizada solo, $\theta$.
		
		\begin{align*}
			\vec{r} = \left( x \;,\; y\right)  = \left(  l\sin(\theta) \;,\; -l\cos(\theta) \right) 
		\end{align*}
		
		Veamos la velocidad de la masa. Como \comillas{S} es inercial, solamente tenemos que derivarlo sin tener mas consideraciones.
		
		\begin{align*}
			\dot{\vec{r}} = \left( l\cos(\theta)\dot{\theta} \;,\; l\sin(\theta)\dot{\theta} \right) 
		\end{align*}
		
		Veamos la energía cinética \comillas{T}
		
		\begin{align*}
			T = \frac{m}{2} \left| \dot{\vec{r}} \right|^2  = \frac{m}{2}\sqrt{(l^2\cos^2(\theta)\dot{\theta}^2 + l^2\sin^2(\theta)\dot{\theta}^2)^2} = \frac{m}{2}l^2\dot{\theta}^2
		\end{align*}
		
		Veamos la energía potencial \comillas{V}, en donde mi origen de potencial $V = 0$ está justo situado en el techo., la masa esta situado a una altura de $h = y = l\sin(\theta) \;,\; -l\cos(\theta)$
		
		\begin{align*}
			V = mgy = mg\left( -l\cos(\theta) \right) 
		\end{align*}
		
		Veamos la energía $E = T + V$
		\begin{align}
			E = T + V = \frac{m}{2}l^2\dot{\theta}^2 - mg\left( l\cos(\theta) \right) 
			\label{eq: E}
		\end{align}
		
		Como no consideramos ninguna fuerza disipativa, la energía se conserva, y por ende la derivada temporal de la energía es 0.
		
		\begin{align*}
			\frac{dE}{dt} = 0
			\\
			0 = ml^2\dot{\theta}\ddot{\theta} + mgl\sin(\theta)\dot{\theta}
		\end{align*}
		
		Imponemos que $\dot{\theta} \neq 0$, pues si no es 0, entonces no hay movimiento, así podemos dividir entre $\dot{\theta}$:
		
		\begin{align*}
			0 = ml^2\ddot{\theta} + mgl\sin(\theta)
		\end{align*}
		
		\begin{align}
			\boxed{\ddot{\theta} + \frac{g}{l}\sin(\theta) = 0}
		\end{align}	
		
		Esta expresión será la que definirá el movimiento de nuestra masa. En el código se utilizó de la siguiente manera
		\begin{align*}
			\begin{cases}
				\frac{d\theta}{dt} = \dot{\theta}
				\\
				\frac{d^2\theta}{dt^2} = -\frac{g}{l}\sin(\theta)
			\end{cases}
		\end{align*}
		
		Si ponemos que $y_o = \theta ; y_1 = \dot{\theta} ; \dfrac{dy_1}{dt} = \ddot{\theta}$, entonces:
		
		\begin{equation}
			\boxed{
			\begin{cases}
				\frac{dy_0}{dt} = y_1 
				\\
				\frac{dy_1}{dt} = -\frac{g}{L} \sin(y_0)
			\end{cases}
			}	
		\end{equation}
		
		Esta es la que se implementará tal cual en el código
		
		
		\section{Periodo}
		En el programa, será necesario tomar un tiempo inicial y una final. Para el final decidí tomar 2 veces el periodo. De esta manera, la animación volver a su posición inicial y al actualizarse para volver a empezar, habría un corte en la animación.\\
		\\
		Para ello no podemos tomar que el periodo sea $2\pi\sqrt{\frac{l}{g}}$, pues esto equivaldría a uno en el que aproximamos par ángulos pequeños. En su lugar tomaré uno corregido para ángulos grandes \cite{marion2014dinamica}:
		
		En el punto de máxima amplitud $\theta_0$, la velocidad es nula ($\dot{\theta} = 0$), por lo que:
		\begin{equation}
			E_0 = -mgL\cos\theta_0
			\label{eq: E_0}
		\end{equation}
		
		Como la energía se conserva, podemos igualar \hyperref[eq: E_0]{(\ref{eq: E_0})} con \hyperref[eq: E]{(\ref{eq: E})}:
		\begin{equation*}
			\frac{1}{2} m L^2 \left( \frac{d\theta}{dt} \right)^2 - mgL\cos\theta = -mgL\cos\theta_0
		\end{equation*}
		
		Despejamos la velocidad angular $\frac{d\theta}{dt}$:
		\begin{gather*}
			\left( \frac{d\theta}{dt} \right)^2 = \frac{2g}{L} (\cos\theta - \cos\theta_0)
			\\
			\frac{d\theta}{dt} = \sqrt{\frac{2g}{L}} \sqrt{\cos\theta - \cos\theta_0}
		\end{gather*}
		
		El periodo total $T$ es cuatro veces el tiempo invertido en desplazarse desde $\theta = 0$ hasta $\theta = \theta_0$:
		\begin{equation*}
			T = 4 \int_{0}^{\theta_0} \sqrt{\frac{L}{2g}} \frac{d\theta}{\sqrt{\cos\theta - \cos\theta_0}}
		\end{equation*}
		
		Utilizando la identidad trigonométrica $\cos\alpha = 1 - 2\sin^2(\alpha/2)$:
		\begin{equation*}
			\cos\theta - \cos\theta_0 = 2 \left( \sin^2\frac{\theta_0}{2} - \sin^2\frac{\theta}{2} \right)
		\end{equation*}
		
		Sustituyendo en la integral:
		\begin{equation*}
			T = 4\sqrt{\frac{L}{g}} \int_{0}^{\theta_0} \frac{d\theta}{2\sqrt{\sin^2\frac{\theta_0}{2} - \sin^2\frac{\theta}{2}}}
		\end{equation*}
		
		La solución exacta de esta integral es una integral elíptica completa de primera especie. Su expansión en serie de potencias es:
		\begin{equation*}
			T = 2\pi\sqrt{\frac{L}{g}} \left( 1 + \frac{1}{4}\sin^2\frac{\theta_0}{2} + \frac{9}{64}\sin^4\frac{\theta_0}{2} + \dots \right)
		\end{equation*}
		
		Para ángulos pequeños, aplicamos la aproximación $\sin\frac{\theta_0}{2} \approx \frac{\theta_0}{2}$ (en radianes). Elevando al cuadrado:
		\begin{equation}
			\sin^2\frac{\theta_0}{2} \approx \left( \frac{\theta_0}{2} \right)^2 = \frac{\theta_0^2}{4}
		\end{equation}
		
		Sustituyendo este término en la expansión:
		\begin{equation}
			T \approx 2\pi\sqrt{\frac{L}{g}} \left( 1 + \frac{1}{4} \cdot \frac{\theta_0^2}{4} \right) = 2\pi\sqrt{\frac{L}{g}} \left( 1 + \frac{1}{16}\theta_0^2 \right)
		\end{equation}
		
		Finalmente, para $10$ ciclos y convirtiendo $\theta_0$ de grados a radianes:
		\begin{equation}
			\boxed{T_{final} = 10 \cdot 2\pi\sqrt{\frac{L}{g}} \left( 1 + \frac{1}{16} \left( \theta_0 \cdot \frac{\pi}{180} \right)^2 \right)}
		\end{equation}
		
		\renewcommand{\refname}{Bibliografía}
		\bibliographystyle{plain}  % O puedes usar otro estilo: ieeetr, apa, etc.
		\bibliography{Bibliografia}
\end{document}
